
\subsection*{\SGAInlineRefTarget{sga3:viib:subsection:4-4}4.4 Theorem}
\phantomsection
\label{sga3:VIIB:4.4}

\noindent\textbf{Theorem 4.4.}
{\itshape Let \(k\) be a pseudocompact ring of characteristic
\(p>0\), and let \(G\) be a formal \(k\)-group. The following
assertions are equivalent:
\begin{enumerate}
\item[\textup{(i)}] There exists a \(p\)-Lie algebra \(\mathbf L\),
flat over \(\mathcal O_k\), such that \(G\) is isomorphic to
\(\mathscr G_p(\mathbf L)\). In this case,
\(\mathbf L=\operatorname{Lie}(G)\) by \SGARefLink{sga3:VIIB:4.2.2}{4.2.2}.

\item[\textup{(ii)}] There exists a projective pseudocompact
\(k\)-module \(\omega\) such that the affine algebra of \(G\) is
isomorphic to the quotient of \(k[[\omega]]\) by the closed ideal
generated by the elements \(x^p\), \(x\in\omega\). In this case,
\(\omega\simeq\omega_{G/k}\).

\item[\textup{\SGAInlineRefTarget{sga3:viib:subsection:4-4:item:iii}(iii)}] The group \(G\) has height at most \(1\), and
\(\omega_{G/k}\) is a projective pseudocompact \(k\)-module.

\item[\textup{(iv)}] The group \(G\) has height at most \(1\), and
is topologically flat over \(k\).
\end{enumerate}}

\noindent\textit{Proof of \textup{\SGARefLink{sga3:vii:paragraph:5-1:equation:i}{(i)}} \(\boldsymbol{\Rightarrow}\)
\textup{\SGARefLink{sga3:vii:paragraph:5-1:equation:ii}{(ii)}}.}
Set
\[
  \omega=\Gamma^*(\mathbf L)
\]
\textup{(cf.\ \SGARefLink{sga3:VIIB:1.2.3.D}{1.2.3.D})}, and let \(A\) be the quotient of
\(k[[\omega]]\) by the closed ideal generated by the elements
\(x^p\), \(x\in\omega\). Every continuous morphism
\[
  h:A\longrightarrow C,
  \qquad C\in\operatorname{Alf}/k,
\]
is determined by its restriction \(h'\) to \(\omega\), and this
restriction maps \(\omega\) into \(\sqrt[p]{0_C}\). Thus there is a
canonical bijection from
\[
  \operatorname{Hom}_{\operatorname{Alp}/k}(A,C)
\]
onto the set
\[
  \operatorname{Hom}_{c}\bigl(\omega,\sqrt[p]{0_C}\bigr)
\]
of continuous \(k\)-linear maps from \(\omega\) into
\(\sqrt[p]{0_C}\). The latter set is canonically isomorphic to
\[
  \mathbf L(C)\otimes_C\sqrt[p]{0_C}
\]
\textup{(see the proof of \SGARefLink{sga3:VIIB:3.3:cartier-theorem}{3.3})}. The implication
\textup{\SGARefLink{sga3:vii:paragraph:5-1:equation:i}{(i)}} \(\Rightarrow\) \textup{\SGARefLink{sga3:vii:paragraph:5-1:equation:ii}{(ii)}} therefore follows from
the functorial bijection
\[
  \mathbf L(C)\otimes_C\sqrt[p]{0_C}
  \xrightarrow{\ \sim\ }
  \mathscr G_p(\mathbf L)(C)
\]
established in \SGARefLink{sga3:VIIA:4.3.1}{4.3.1}.

For the remaining implications, consult the analogous proofs of
\SGARefLink{sga3:VIIB:3.3:cartier-theorem}{Theorems~3.3} and VIIA~\SGARefLink{sga3:VIIA:7.4.1}{7.4.1}.

\medskip
\noindent\textbf{Remark 4.4.A.}
\phantomsection
\label{sga3:VIIB:4.4.A}
Let \(G\) be an infinitesimal formal \(k\)-group, with affine
algebra \(A\), such that
\[
  \omega_{G/k}=I_A/I_A^2
\]
is a projective pseudocompact \(k\)-module. Then the projection
\[
  I_A\longrightarrow\omega_{G/k}
\]
has a section
\[
  \tau:\omega_{G/k}\longrightarrow I_A,
\]
and \(\tau\) induces a surjective continuous morphism of algebras
\[
  \psi:k[[\omega_{G/k}]]\longrightarrow A;
\]
compare the proof of the implication
\textup{\SGARefLink{sga3:viib:subsection:4-4:item:iii}{(iii)}} \(\Rightarrow\) \textup{\SGARefLink{sga3:vii:paragraph:5-1:equation:i}{(i)}} in \SGARefLink{sga3:VIIB:3.3:cartier-theorem}{3.3}.\footnote{%
\label{sga3:VIIB:nde:139}The editors added this remark, which is
used in \SGARefLink{sga3:VIIB:4.4.2}{4.4.2}.}

\subsubsection*{4.4.1 Corollary}
\phantomsection
\label{sga3:VIIB:4.4.1}

\noindent\textbf{Corollary 4.4.1.}
{\itshape If \(k\) is a field of characteristic \(p>0\), the functor
\[
  \mathbf L\longmapsto\mathscr G_p(\mathbf L)
\]
is an equivalence from the category of \(p\)-Lie algebras over
\(k\) to the category of formal \(k\)-groups of height at most
\(1\).}

Indeed, as \(G\) ranges over the formal groups of height at most
\(1\), \SGARefLink{sga3:viib:subsection:4-4}{Theorem~4.4} identifies the functor
\[
  G\longmapsto\mathscr G_p\bigl(\operatorname{Lie}G\bigr)
\]
with the identity functor. Likewise, \SGARefLink{sga3:VIIB:4.2.2}{4.2.2} identifies the functor
\[
  \mathbf L\longmapsto
  \operatorname{Lie}\mathscr G_p(\mathbf L)
\]
with the identity functor \textup{(see also VIIA~\SGARefLink{sga3:VIIA:5.5.3}{5.5.3})}.\footnote{%
\label{sga3:VIIB:nde:140}If \(k\) is an Artinian ring of
characteristic \(p>0\), the same proof gives an equivalence between
the category of \(p\)-Lie algebras flat over \(k\) and the category
of formal \(k\)-groups of height at most \(1\) that are
topologically flat over \(k\).}
